\documentclass[11pt]{article}

\usepackage[T1]{fontenc}
\usepackage{lmodern}
\usepackage[a4paper,margin=29mm]{geometry}
\usepackage{amsmath,amssymb,amsthm,mathtools}
\usepackage{microtype}
\usepackage{enumitem}
\usepackage{xcolor}
\usepackage[colorlinks=true,linkcolor=blue!45!black,citecolor=blue!45!black,
  urlcolor=blue!55!black]{hyperref}
\setlength{\emergencystretch}{2em}

\newtheorem{theorem}{Theorem}[section]
\newtheorem{proposition}[theorem]{Proposition}
\newtheorem{lemma}[theorem]{Lemma}
\newtheorem{corollary}[theorem]{Corollary}
\newtheorem{definition}[theorem]{Definition}
\newtheorem{remark}[theorem]{Remark}

\newcommand{\C}{\mathbb C}
\newcommand{\R}{\mathbb R}
\newcommand{\cS}{\mathcal S}
\newcommand{\cR}{\mathcal R}
\newcommand{\cZ}{\mathcal Z}
\newcommand{\supp}{\operatorname{supp}}
\newcommand{\ran}{\operatorname{ran}}
\newcommand{\Span}{\operatorname{span}}

\title{No Positive-Definite Metric on the Schwartz Source Core\\
of the Hedenmalm Theta Pair}
\author{Michael Haufschild}
\date{}

\begin{document}

\maketitle

\begin{abstract}
Hedenmalm recently represented the zeros of the centered Riemann function
\(\Xi(z)=\xi(\frac12+iz)\) by a theta-driven differential pencil and proposed a
self-adjointness problem for the pair \((LD,L)\).  We determine the metric geometry of this
proposal on its natural translation-stable source core.

Let \(k\) be the normalized logarithmic theta source, so that \(\widehat{k}=\Xi\), let
\(D=-i\,d/dx\), and let \(Q\) be a jointly continuous nonnegative sesquilinear form on complex
Schwartz space.  If \(D\) is symmetric for \(Q\) and \(Q(k,k)=0\), then \(Q\) has a unique
representation
\[
 Q(f,g)=\int_{\R}\widehat f(t)\overline{\widehat g(t)}\,d\mu(t),
\]
where \(\mu\) is a positive measure of temperate growth supported on the real zero set of
\(\Xi\).  Consequently \(\mu\) is purely atomic.  Conversely, every such divisor-supported
measure gives a form with these properties.

The classification has a sharp algebraic and topological consequence.  If \(Q\neq0\), the image in ordinary
\(L^2(\R)\) of its radical is dense, proper, and nonclosed.  Moreover, for every source vector
\(f\) there is a sequence \(z_n\to0\) in \(L^2\) with
\[
 Q(z_m-z_n,z_m-z_n)=0,\qquad Q(z_n,z_n)=Q(f,f).
\]
Thus sequential \(L^2\)-closability is equivalent to \(Q=0\).  In particular, no nonzero
source-core metric for the Hedenmalm pair can be the restriction of a closed nonnegative
quadratic form on ordinary \(L^2\).  More strongly, if \(Q\) is pulled back from a
positive-definite form on the full algebraic range \(L\mathcal S(\R)\), its radical must equal
\(\ker L=\C k\), contradicting the dense-radical theorem.  Hence no such positive-definite
source-core metric exists, even before completion.

For the Fourier eigenprofile attached to a zero \(\alpha\), the classified atomic polarization is
zero when \(\alpha\notin\R\), while at a real zero it equals
\(\mu(\{\alpha\})|\Xi'(\alpha)|^2\).  Hence conservativity on all theta eigenprofiles forces
both the Riemann hypothesis and simplicity.  The result does not assert RH; it proves that the
ordinary Hilbert-space completion of this new spectral program collapses and that every
nonzero source-core alternative is necessarily a degenerate, singular, divisor-supported
Fourier rigging.
\end{abstract}

\section{Introduction}

The Hilbert--Pólya principle seeks a self-adjoint spectral mechanism whose eigenvalues encode the
nontrivial zeros of the Riemann zeta function.  The principle is compelling because
self-adjointness would force the spectral parameters to be real, but a spectral
parametrization alone is not enough: the decisive issue is the construction of a positive
Hilbert geometry \cite{Lagarias2006,Connes2026}.  This distinction is especially visible in the
recent theta spectral problem of Hedenmalm \cite{Hedenmalm2026}.

Write
\[
 \xi(s)=\frac12s(s-1)\pi^{-s/2}\Gamma(s/2)\zeta(s),
 \qquad
 \Xi(z)=\xi\!\left(\frac12+iz\right).
 \tag{1.1}
\]
Hedenmalm constructs first-order differential operators \(D\) and \(L\) on the positive
half-line and shows that the boundary-value problem
\[
 LDu+\alpha Lu=0
 \tag{1.2}
\]
has a nonzero solution precisely when \(\Xi(\alpha)=0\), including complex \(\alpha\).
He then asks for a Hilbert form that makes the pair \((LD,L)\) self-adjoint.  If that form is a
genuine Hilbert norm on the required image space, the standard symmetry calculation excludes
nonreal \(\alpha\) and proves RH.

The present paper answers the metric-classification problem when the proposed form extends from
the eigenfunction span to the natural Schwartz source core.  This regularity boundary is both
strong enough for the differential pair to act and weak enough to include distributional
Fourier-sampling metrics.  The answer is rigid:
\[
\boxed{
\begin{gathered}
\text{positive source-core pair symmetry}\\
\Longrightarrow
\text{positive Fourier measure on the real \(\Xi\)-divisor}\\
\Longrightarrow
\begin{matrix}
\text{dense nonclosed \(L^2\)-radical},\\
\text{no positive-definite metric on \(L\cS(\R)\)}.
\end{matrix}
\end{gathered}}
\tag{1.3}
\]
Thus the metric cannot be repaired within ordinary \(L^2\), even by passing from bounded
operators to closed unbounded quadratic forms.  The only nonzero metrics in the classified
category are singular sampling riggings.

\subsection{The source-core metric problem}

After the logarithmic change of variables, the theta pair takes the form
\[
 D=-i\frac{d}{dx},\qquad L=kDk^{-1},
 \tag{1.4}
\]
where \(k\) is a positive Schwartz function satisfying
\[
 Lk=0,\qquad \widehat k=\Xi.
 \tag{1.5}
\]
Suppose a candidate form \(q_2\) is defined on the image of \(L\).  Its pullback is
\[
 Q(u,v)=q_2(Lu,Lv).
 \tag{1.6}
\]
The pair-symmetry identity is exactly
\[
 Q(Du,v)=Q(u,Dv),
 \tag{1.7}
\]
and \(Lk=0\) gives \(Q(k,k)=0\).  We study all jointly continuous nonnegative
sesquilinear forms on \(\cS(\R)\) satisfying these two identities.  No zero location, RH,
simplicity, or positivity of a zeta-derived kernel is assumed.

Here \(L\cS(\R)\) means only the algebraic image of the Schwartz source inside the ambient
smooth target class.  We do not assume that \(L\) maps \(\cS(\R)\) into itself or into ordinary
\(L^2\); its logarithmic coefficient is unbounded.  The regularity hypothesis is solely that the
pullback \(Q\), not the target form \(q_2\), is jointly continuous in the Schwartz topology.
Thus the theorem tests a source-core extension without smuggling in a target-space completion.

\subsection{Main theorem}

Use the Fourier convention
\[
 \widehat f(t)=\int_{\R}f(x)e^{itx}\,dx.
 \tag{1.8}
\]
A positive Borel measure \(\mu\) has \emph{temperate growth} if
\[
 \int_{\R}(1+t^2)^{-N}\,d\mu(t)<\infty
 \tag{1.9}
\]
for some \(N\ge0\).  Let
\[
 \cZ_{\R}=\{t\in\R:\Xi(t)=0\}.
 \tag{1.10}
\]

\begin{theorem}[Divisor metric and dense-radical law]\label{thm:main}
Let \(Q:\cS(\R)\times\cS(\R)\to\C\) be jointly continuous, complex-linear in the first
argument, conjugate-linear in the second, and nonnegative.  Assume
\[
 Q(Df,g)=Q(f,Dg),\qquad Q(k,k)=0.
 \tag{1.11}
\]
Then:
\begin{enumerate}[label=\textup{(\roman*)}]
\item There is a unique positive measure \(\mu\) of temperate growth such that
\[
 Q(f,g)=\int_{\R}\widehat f(t)\overline{\widehat g(t)}\,d\mu(t).
 \tag{1.12}
\]
\item Its support is contained in \(\cZ_{\R}\).  Hence
\[
 \mu=\sum_{\gamma\in\cZ_{\R}}c_\gamma\delta_\gamma,
 \qquad c_\gamma\ge0,
 \tag{1.13}
\]
with the weights satisfying the temperate-growth condition.
\item Conversely, every positive temperate measure of the form \textup{(1.13)} defines a
jointly continuous nonnegative form satisfying \textup{(1.11)}.
\item Let
\[
 \cR_Q=\{f\in\cS(\R):Q(f,g)=0\ \text{for every }g\in\cS(\R)\}.
 \tag{1.14}
\]
If \(Q\neq0\), the image of \(\cR_Q\) in \(L^2(\R)\) is dense, proper, and nonclosed.
\item If \(Q\neq0\), then
\[
 \C k\subsetneq\cR_Q.
 \tag{1.15}
\]
Thus the induced form on \(L\cS(\R)\cong\cS(\R)/\C k\) is never positive definite.
\item For every \(f\in\cS(\R)\) there is a sequence \(z_n\in\cS(\R)\) such that
\[
\begin{split}
 &\|z_n\|_2\longrightarrow0,\\
 &Q(z_m-z_n,z_m-z_n)=0\quad(m,n\ge1),\\
 &Q(z_n,z_n)=Q(f,f)\quad(n\ge1).
\end{split}
\tag{1.16}
\]
Consequently \(Q\) is sequentially closable through \(\cS(\R)\hookrightarrow L^2(\R)\) if
and only if \(Q=0\).
\end{enumerate}
\end{theorem}

The strict containment in (1.15) is the algebraic no-go.  The fixed-energy sequence in (1.16) is
the topological no-go and is stronger than failure of bounded representability.  It
rules out every closed nonnegative \(L^2\)-form whose restriction is \(Q\).

\begin{corollary}[Source-core no-go for the theta pair]\label{cor:no-go}
Let \(q_2\) be a nonnegative form on a space containing \(L\cS(\R)\), and suppose its pullback
\(Q(f,g)=q_2(Lf,Lg)\) is a nonzero jointly continuous form on \(\cS(\R)\).  If the Hedenmalm pair
is symmetric on this core, then:
\begin{enumerate}[label=\textup{(\alph*)}]
\item \(Q\) is represented by a positive atomic measure on \(\cZ_{\R}\);
\item the \(L^2\)-image of its radical is not closed;
\item \(Q\) has no jointly continuous extension to \(L^2(\R)\);
\item \(Q\) is not the restriction of a closed nonnegative quadratic form on \(L^2(\R)\).
\end{enumerate}
\end{corollary}

\begin{corollary}[No positive-definite source-core metric]\label{cor:no-hilbert}
There is no positive-definite sesquilinear form \(q_2\) on the algebraic range
\(L\cS(\R)\) for which
\[
 Q(f,g)=q_2(Lf,Lg)
 \tag{1.17}
\]
is jointly continuous on \(\cS(\R)\times\cS(\R)\) and makes the Hedenmalm pair symmetric on
that source core.
\end{corollary}

\begin{proof}
If such a form existed, positivity would give
\[
 f\in\cR_Q
 \quad\Longleftrightarrow\quad
 q_2(Lf,Lf)=0
 \quad\Longleftrightarrow\quad
 Lf=0.
\]
Thus \(\cR_Q=\ker L=\C k\).  The form \(Q\) is nonzero: \(xk\in\cS(\R)\),
\(L(xk)=-ik\ne0\), and hence \(Q(xk,xk)>0\).  Theorem \ref{thm:main}\textup{(iv)} now says that
the \(L^2\)-image of \(\C k\) is dense.  This is impossible because a nonzero one-dimensional
subspace of \(L^2(\R)\) is closed and proper.
\end{proof}

\subsection{What is and is not excluded}

Corollary \ref{cor:no-hilbert} excludes a positive-definite pair metric on the full algebraic
range \(L\cS(\R)\) whenever its pullback has the natural Schwartz continuity.  Corollary
\ref{cor:no-go} separately excludes every ordinary-\(L^2\), translation-stable
Schwartz-source completion of a nonzero semidefinite metric.  Neither result concerns a form defined only on the
finite linear span of eigenfunctions and refusing every continuous source-core extension.  Such
a target-domain construction lies outside the hypotheses.  This boundary is exact rather than
cosmetic: Fourier evaluation at a point is continuous on Schwartz space and discontinuous on
ordinary \(L^2\), which is precisely the singularity forced by Theorem \ref{thm:main}.

The novelty is not the Bochner--Schwartz theorem or Wiener's cyclicity theorem separately.
It is their source-specific combination with (1.5)--(1.7), which converts Hedenmalm's metric
problem into the atomic classification (1.13), strict source degeneracy (1.15), and the maximal
\(L^2\)-topology obstruction (1.16).  Section \ref{sec:literature} gives the comparison with the
direct parent.

\section{The normalized theta source}\label{sec:source}

This section fixes normalization without relying on a zero-dependent construction.  Put
\[
 h(t)=2\pi(2\pi t^4-3t^2)e^{-\pi t^2},\qquad t>0,
 \tag{2.1}
\]
and define
\[
k(x)=e^{x/2}\sum_{n=1}^{\infty}h(ne^x).
\tag{2.2}
\]
Termwise comparison with Hedenmalm's theta derivative gives
\[
 k(x)=2\,\mathop{\mathrm{Var}\Theta}_{00}(i e^{2x}),
\]
so this is precisely his logarithmic source with the scalar normalization that makes
\(\widehat k=\Xi\).  Multiplying \(k\) by a nonzero constant does not change \(L=kDk^{-1}\).

\begin{proposition}[Theta--Fourier identity]\label{prop:theta-fourier}
The function \(k\) is nonzero, even, smooth, and belongs to \(\cS(\R)\).  It is strictly positive,
and its Fourier transform is
\[
 \widehat k(z)=\Xi(z)\qquad(z\in\C).
 \tag{2.3}
\]
\end{proposition}

\begin{proof}
The Mellin transform of \(h\) is
\[
\begin{aligned}
 \int_0^\infty h(t)t^{s-1}\,dt
 &=\frac12s(s-1)\pi^{-s/2}\Gamma(s/2).
\end{aligned}
\tag{2.4}
\]
For \(\Re s>1\), absolute convergence and the substitution \(u=ne^x\) give
\[
\begin{aligned}
 \int_{\R}k(x)e^{i z x}\,dx
 &=\sum_{n\ge1}\int_0^\infty h(nt)t^{s}\,\frac{dt}{t}\\
 &=\left(\sum_{n\ge1}n^{-s}\right)
   \int_0^\infty h(u)u^{s-1}\,du\\
 &=\xi(s),\qquad s=\frac12+iz.
\end{aligned}
\tag{2.5}
\]
The Jacobi inversion formula applied to the theta series in (2.2) makes \(k\) even and gives
double-exponential decay, with the same decay for every derivative.  It also gives strict
positivity; these are the theta estimates used in the construction of
\cite{Hedenmalm2026}.  Thus \(k\in\cS(\R)\), and the left side of (2.5) is entire in \(z\).
Analytic continuation proves (2.3) on \(\C\).
\end{proof}

Define
\[
 \Phi(x)=-\log k(x),\qquad
 D=-i\partial_x,\qquad
 L=-i(\partial_x+\Phi'(x))=kDk^{-1}.
 \tag{2.6}
\]
Then
\[
 Lk=0,\qquad \ker L=\C k
 \tag{2.7}
\]
on the smooth source class.

\begin{remark}
The scalar normalization in (2.1)--(2.2) is chosen solely to make (2.3) exact.  Since
\(\Phi\) changes only by an additive constant under a scalar change of \(k\), the differential
operator \(L\), its kernel line, and every result below are normalization invariant.
\end{remark}

\section{From pair symmetry to translation invariance}\label{sec:translation}

Let \(T_a f(x)=f(x+a)\).  Joint continuity of \(Q\) is taken in the usual Schwartz topology.
Nonnegativity means
\[
 Q(f,f)\in[0,\infty)\qquad(f\in\cS(\R)).
\tag{3.1}
\]
It implies Hermitian symmetry and the Cauchy--Schwarz inequality for the associated seminorm.

\begin{lemma}[Theta diagonal nullity is radical nullity]\label{lem:null-radical}
If \(Q\) is nonnegative and \(Q(k,k)=0\), then
\[
 Q(k,g)=Q(g,k)=0\qquad(g\in\cS(\R)).
\tag{3.2}
\]
\end{lemma}

\begin{proof}
The Cauchy--Schwarz inequality gives
\[
 |Q(k,g)|^2\le Q(k,k)Q(g,g)=0.
\]
Hermitian symmetry gives the other identity.
\end{proof}

\begin{lemma}[Momentum symmetry]\label{lem:translation}
If
\[
 Q(Df,g)=Q(f,Dg)\qquad(f,g\in\cS(\R)),
\tag{3.3}
\]
then
\[
 Q(T_a f,T_a g)=Q(f,g)\qquad(a\in\R).
\tag{3.4}
\]
\end{lemma}

\begin{proof}
Since \(D=-i\partial_x\), (3.3) is equivalent to
\[
 Q(\partial_x f,g)+Q(f,\partial_x g)=0.
\tag{3.5}
\]
The map \(a\mapsto T_a f\) is differentiable in the Schwartz topology with derivative
\(T_a\partial_x f\).  Joint continuity therefore permits differentiation:
\[
\begin{aligned}
 \frac{d}{da}Q(T_a f,T_a g)
 &=Q(T_a\partial_x f,T_a g)+Q(T_a f,T_a\partial_x g)=0
\end{aligned}
\]
by (3.5).  Hence the matrix coefficient is constant, proving (3.4).
\end{proof}

\begin{proposition}[Pair reduction]\label{prop:pair-reduction}
Let \(q_2\) be a nonnegative sesquilinear form on a space containing \(L\cS(\R)\), and put
\[
 Q(f,g)=q_2(Lf,Lg).
\tag{3.6}
\]
Then \(Q(k,k)=0\).  Moreover,
\[
 q_2(LDf,Lg)=q_2(Lf,LDg)
\tag{3.7}
\]
is equivalent to momentum symmetry (3.3) for \(Q\).
\end{proposition}

\begin{proof}
The first assertion follows from \(Lk=0\).  Substituting (3.6) into (3.3) gives (3.7) exactly.
\end{proof}

\section{Classification by the real \(\Xi\)-divisor}\label{sec:classification}

We use the Bochner--Schwartz theorem in the following standard form
\cite[Chap.~VII, \S9]{Schwartz1966}: a positive-definite tempered distribution on \(\R\) is the
Fourier transform of a unique positive measure of temperate growth.

\begin{proposition}[Positive translation-invariant Schwartz forms]\label{prop:BS}
Let \(Q\) be a jointly continuous nonnegative sesquilinear form on \(\cS(\R)\).  If
\[
 Q(T_a f,T_a g)=Q(f,g)
\tag{4.1}
\]
for all \(a,f,g\), then there is a unique positive measure \(\mu\) of temperate growth such that
\[
 Q(f,g)=\int_{\R}\widehat f(t)\overline{\widehat g(t)}\,d\mu(t).
\tag{4.2}
\]
Conversely, every positive measure of temperate growth defines a jointly continuous
nonnegative translation-invariant form by (4.2).
\end{proposition}

\begin{proof}
By the Schwartz kernel theorem, joint continuity represents \(Q\) by a tempered distribution
on \(\R^2\).  Invariance under simultaneous translation makes that kernel depend only on the
difference variable.  Nonnegativity makes the resulting tempered distribution positive
definite.  The Bochner--Schwartz theorem gives a positive temperate measure \(\mu\), and Fourier
polarization gives (4.2).

Uniqueness follows because Fourier transforms of Schwartz functions contain
\(C_c^\infty(\R)\): if two measures give the same form, compactly supported smooth spectral
bumps and polarization make their integrals agree on \(C_c^\infty(\R)\), hence the measures
agree.  Conversely, (1.9) and rapid decay of \(\widehat f,\widehat g\) imply both convergence
and joint Schwartz continuity of (4.2).  Positivity and translation invariance are immediate.
\end{proof}

\begin{theorem}[Divisor-support classification]\label{thm:divisor}
Under the hypotheses of Theorem \ref{thm:main}, the measure in
Proposition \ref{prop:BS} satisfies
\[
 \supp\mu\subset\cZ_{\R}.
\tag{4.3}
\]
It is therefore purely atomic as in (1.13).  Conversely, every positive temperate measure
supported on \(\cZ_{\R}\) satisfies \(Q(k,k)=0\).
\end{theorem}

\begin{proof}
Lemmas \ref{lem:translation} and Proposition \ref{prop:BS} give the unique measure.  By
Proposition \ref{prop:theta-fourier},
\[
 0=Q(k,k)=\int_{\R}|\Xi(t)|^2\,d\mu(t).
\tag{4.4}
\]
The integrand is continuous and nonnegative.  Hence it vanishes on \(\supp\mu\), proving
(4.3).

The real zeros of a nonzero entire function are discrete and have no finite accumulation point.
A positive measure supported on this set is the sum of its point masses, which gives (1.13).
The temperate-growth condition is inherited from \(\mu\).  Conversely, if (4.3) holds, the
integral in (4.4) is zero.
\end{proof}

\begin{corollary}[Descent through the theta constraint]\label{cor:descent}
Every form in Theorem \ref{thm:divisor} has \(k\) in both radicals and descends to a
well-defined nonnegative form on the algebraic range of \(L\) by
\[
 q_2(Lf,Lg)=Q(f,g).
\tag{4.5}
\]
\end{corollary}

\begin{proof}
Radical nullity follows from Lemma \ref{lem:null-radical}.  If \(Lf=Lf'\), then
\(f-f'\in\ker L=\C k\), so (4.5) is independent of the chosen preimages.
\end{proof}

The theorem is a complete classification: positivity and source-core pair symmetry leave only
the weights \(c_\gamma\).  No absolutely continuous part, singular continuous part, or mass away
from real zeta zeros survives.

\section{Cyclicity and the dense-radical law}\label{sec:cyclicity}

The divisor classification already shows that nonzero metrics are singular.  A second argument,
using only source cyclicity, strengthens singularity to nonclosability.

\begin{proposition}[Theta cyclicity]\label{prop:cyclic}
The normalized theta source is cyclic for real translations in ordinary \(L^2(\R)\):
\[
 \overline{\Span\{T_a k:a\in\R\}}^{\,L^2}=L^2(\R).
\tag{5.1}
\]
\end{proposition}

\begin{proof}
The real zero set of the nonzero entire function \(\Xi\) is countable and therefore
Lebesgue-null.  Proposition \ref{prop:theta-fourier} gives
\[
 \widehat k(t)\ne0\quad\text{for almost every }t\in\R.
\tag{5.2}
\]
Wiener's \(L^2\) cyclicity theorem \cite{Wiener1932} now gives (5.1).

For completeness, suppose \(h\in L^2(\R)\) is orthogonal to every translated source.  Under
Plancherel, translation becomes multiplication by a character, so the finite complex measure
\[
 d\nu(t)=h(t)\overline{\Xi(t)}\,dt
\tag{5.3}
\]
has identically zero Fourier transform.  Uniqueness of the Fourier transform of finite measures
gives \(\nu=0\).  Equation (5.2) then gives \(h=0\).  The orthogonal complement of the orbit span
is zero, proving (5.1).
\end{proof}

\begin{lemma}[The translation orbit lies in the radical]\label{lem:orbit-radical}
Under the hypotheses of Theorem \ref{thm:main},
\[
 T_a k\in\cR_Q\qquad(a\in\R).
\tag{5.4}
\]
\end{lemma}

\begin{proof}
Given \(g\in\cS(\R)\), use Lemmas \ref{lem:null-radical} and \ref{lem:translation}:
\[
 Q(T_a k,g)
 =Q(T_a k,T_aT_{-a}g)
 =Q(k,T_{-a}g)=0.
\]
\end{proof}

\begin{theorem}[Dense proper radical]\label{thm:dense-radical}
If \(Q\neq0\), the image of \(\cR_Q\) under the inclusion
\(\cS(\R)\hookrightarrow L^2(\R)\) is dense, proper, and nonclosed.
\end{theorem}

\begin{proof}
Lemma \ref{lem:orbit-radical} and Proposition \ref{prop:cyclic} make the image dense.  If it were
all of \(L^2\), then every \(f\in\cS(\R)\) would have the same \(L^2\)-class as some
\(r\in\cR_Q\).  Equality almost everywhere of two continuous functions gives \(f=r\), so
\(f\in\cR_Q\).  This would make \(Q=0\), contrary to the hypothesis.  Thus the image is proper.
A dense proper linear subspace is not closed.
\end{proof}

\begin{theorem}[Fixed-energy null sequence]\label{thm:null-sequence}
For each \(f\in\cS(\R)\), there is a sequence satisfying (1.16).
\end{theorem}

\begin{proof}
By density in Theorem \ref{thm:dense-radical}, choose \(r_n\in\cR_Q\) with
\[
 \|f-r_n\|_2<\frac1n.
\tag{5.5}
\]
Set \(z_n=f-r_n\).  Then \(z_n\to0\) in \(L^2\).  Because \(\cR_Q\) is a linear subspace and
lies in both radicals,
\[
\begin{aligned}
 Q(z_n,z_n)&=Q(f-r_n,f-r_n)=Q(f,f),\\
 Q(z_m-z_n,z_m-z_n)&=Q(r_n-r_m,r_n-r_m)=0.
\end{aligned}
\]
This is (1.16).
\end{proof}

\begin{corollary}[Strict source degeneracy]\label{cor:strict-degeneracy}
If \(Q\neq0\), then \(\C k\subsetneq\cR_Q\).  Equivalently, the form induced by \(Q\) on
\(\cS(\R)/\ker L\cong L\cS(\R)\) has a nonzero radical.
\end{corollary}

\begin{proof}
Lemma \ref{lem:null-radical} gives \(\C k\subset\cR_Q\).  Equality would make the \(L^2\)-image
of \(\cR_Q\) one-dimensional and therefore closed, contradicting Theorem
\ref{thm:dense-radical}.
\end{proof}

\subsection{Exact closure consequences}

Recall the standard sequential criterion for a nonnegative form \(Q\) on the dense source
\(\cS(\R)\subset L^2(\R)\): it is closable if
\[
\left.
\begin{array}{c}
 f_n\to0\ \text{in }L^2,\\
 Q(f_n-f_m,f_n-f_m)\to0
\end{array}
\right\}
\quad\Longrightarrow\quad
 Q(f_n,f_n)\to0.
\tag{5.6}
\]

\begin{corollary}[Exact \(L^2\)-closure classification]\label{cor:closure}
Under the hypotheses of Theorem \ref{thm:main}, the following are equivalent:
\[
\begin{split}
 &Q\text{ is sequentially closable through }L^2(\R);\\
 &Q=0;\\
 &\mu=0.
\end{split}
\tag{5.7}
\]
\end{corollary}

\begin{proof}
If \(Q\neq0\), positivity and polarization give an \(f\) with \(Q(f,f)>0\).
The sequence from Theorem \ref{thm:null-sequence} violates (5.6), so \(Q\) is not closable.
The zero form is closable.  Finally, (4.2) and uniqueness give \(Q=0\) if and only if
\(\mu=0\).
\end{proof}

\begin{corollary}[No continuous \(L^2\) extension]\label{cor:no-extension}
If a jointly continuous sesquilinear form \(\widetilde Q\) on \(L^2(\R)\) agrees with \(Q\) on
\(\cS(\R)\times\cS(\R)\), then \(Q=\widetilde Q=0\).
\end{corollary}

\begin{proof}
The radical of \(\widetilde Q\) is closed.  It contains the dense translation span of \(k\) by
Lemmas \ref{lem:null-radical}--\ref{lem:orbit-radical} and continuity, so it equals \(L^2\).
Thus \(\widetilde Q=0\), and its restriction is zero.
\end{proof}

Corollaries \ref{cor:closure} and \ref{cor:no-extension} close the two standard ordinary-Hilbert
escapes: neither a bounded metric operator nor a closed unbounded nonnegative form can implement a
nonzero source-core metric.

There is also an obstruction before topology is imposed on the target.  If a form on
\(L\cS(\R)\) were positive definite, its pullback radical would be precisely
\(\ker L=\C k\).  Theorem \ref{thm:dense-radical} forces the same radical to have dense
\(L^2\)-image, whereas \(\C k\) has closed one-dimensional image.  This proves Corollary
\ref{cor:no-hilbert}: every nonzero form surviving the classification is necessarily
semidefinite on \(L\cS(\R)\), with additional null directions beyond the kernel of \(L\).

\section{Exact sampling of theta eigenprofiles}\label{sec:sampling}

The classification also determines what the divisor metric sees on Hedenmalm's
eigenfunction family.  Let \(\alpha\in\C\) satisfy \(\Xi(\alpha)=0\).  The normalized Fourier
profile of the corresponding theta eigenfunction is the entire removable quotient
\[
 F_\alpha(z)=-i\frac{\Xi(z)}{\alpha-z},
\tag{6.1}
\]
with its value at \(z=\alpha\) defined by analytic continuation.  For an atomic measure
\(\mu=\sum_\gamma c_\gamma\delta_\gamma\), extend the spectral polarization to the algebraic
span of these profiles by
\[
 \langle F,G\rangle_\mu
   =\sum_{\gamma\in\cZ_\R}c_\gamma
      F(\gamma)\overline{G(\gamma)}.
\tag{6.2}
\]
Each sum involving finitely many eigenprofiles is well defined: a profile vanishes at every
real divisor point except possibly its own.

\begin{theorem}[Eigenprofile sampling law]\label{thm:eigen-sampling}
Let the polarization have the atomic representation (6.2).
\begin{enumerate}[label=\textup{(\roman*)}]
\item If \(\alpha\notin\R\), then
\[
 \langle F_\alpha,F_\alpha\rangle_\mu=0.
\tag{6.3}
\]
\item If \(\alpha\in\R\), then
\[
 \langle F_\alpha,F_\alpha\rangle_\mu
 =c_\alpha|\Xi'(\alpha)|^2.
\tag{6.4}
\]
\item Profiles attached to distinct real zeros are orthogonal.
\end{enumerate}
\end{theorem}

\begin{proof}
At any real zero \(\gamma\ne\alpha\), equation (6.1) gives
\(F_\alpha(\gamma)=0\).  If \(\alpha\notin\R\), this covers every point of
\(\supp\mu\), proving (6.3).  If \(\alpha\in\R\), removable division gives
\[
 F_\alpha(\alpha)=i\Xi'(\alpha).
\tag{6.5}
\]
Inserting these samples into (6.2) proves (6.4) and the orthogonality assertion.
\end{proof}

\begin{corollary}[Conservativity criterion]\label{cor:conservative}
The atomic polarization is positive definite on the algebraic span of all theta eigenprofiles if and only
if:
\[
\begin{gathered}
\text{every zero of \(\Xi\) is real},\\
\text{every zero of \(\Xi\) is simple},\\
c_\gamma>0\quad\text{at every zero}.
\end{gathered}
\tag{6.6}
\]
In particular, conservativity on all theta eigenprofiles implies RH and simplicity of every
nontrivial zeta zero.
\end{corollary}

\begin{proof}
A nonreal-zero eigenprofile is null by (6.3), and a multiple real zero is null by
\(\Xi'(\alpha)=0\) in (6.4).  If (6.6) holds, then for every finite sum over distinct zeros,
\[
 \left\langle\sum_\alpha a_\alpha F_\alpha,
          \sum_\alpha a_\alpha F_\alpha\right\rangle_\mu
 =\sum_\alpha c_\alpha|a_\alpha|^2|\Xi'(\alpha)|^2,
\tag{6.7}
\]
which vanishes only when every \(a_\alpha=0\).
\end{proof}

This corollary explains the exact limitation of positive semidefinite divisor metrics.
Selecting weights on the real divisor automatically produces pair symmetry, but it annihilates
every hypothetical off-line eigenprofile and every eigenprofile at a multiple real zero.
Positive definiteness is therefore not supplied by the classification; it is equivalent, in this
metric category, to RH together with simplicity and positive mass on every zero fiber.

\section{The resulting Hilbert--Pólya boundary}\label{sec:boundary}

Theorems \ref{thm:divisor}, \ref{thm:dense-radical}, and
\ref{thm:eigen-sampling} give a closed structural answer to the source-core metric problem.

\begin{enumerate}[label=\textup{\arabic*.},leftmargin=2em]
\item \emph{A positive-definite source-core metric is impossible.}
Any positive-definite form on \(L\cS(\R)\) would make the pullback radical exactly
\(\C k\).  Pair symmetry instead forces that radical to be \(L^2\)-dense.  The contradiction
occurs before completing the target space.

\item \emph{Ordinary \(L^2\) is impossible even semidefinitely.}
Every nonzero positive source-core form has a dense nonclosed \(L^2\)-radical and is
nonclosable.  Enlarging a bounded metric to a closed unbounded quadratic form does not evade the
obstruction.

\item \emph{Nonlocality alone is insufficient.}
The representing object is not an arbitrary nonlocal kernel.  It is a unique positive
Fourier measure, and theta nullity forces its complete support onto \(\cZ_{\R}\).

\item \emph{The surviving metric is singular in a precise sense.}
It is a weighted Fourier-evaluation space.  Point evaluation is continuous in the Schwartz
topology but has no continuous or closable realization in the ordinary \(L^2\) topology.

\item \emph{Source-core positivity does not by itself prove RH.}
The positive form sees only real divisor atoms.  Conservativity on the complete eigenprofile
family is the clause that forces RH and simplicity, as stated exactly in
Corollary \ref{cor:conservative}.

\item \emph{A target-domain escape is outside the theorem and visibly different.}
A form defined only on the finite zero-eigenfunction span need not extend continuously to the
translation-stable Schwartz source core.  Such a construction cannot be presented as an ordinary
source \(L^2\)-completion: the latter is excluded by Corollary \ref{cor:no-go}.
\end{enumerate}

Accordingly, the theta-pair architecture has only two mathematically distinct metric regimes:
\[
\begin{gathered}
\text{Schwartz-source extension}
\ \Longrightarrow\ \text{degenerate divisor metric with no \(L^2\)-closure},\\[1mm]
\text{finite target-eigenfunction domain}
\ \Longrightarrow\ \text{no source-core extension assumed}.
\end{gathered}
\tag{7.1}
\]
The first regime is completely classified here.  The second is not an ordinary completion of the
source representation and is not covered by claims made for the first.

\section{Relation to prior work}\label{sec:literature}

Hedenmalm \cite{Hedenmalm2026} supplies the direct parent: the theta boundary-value
characterization, the pair \((LD,L)\), the finite eigenfunction space, and the implication from a
positive self-adjoint pair to RH.  His paper does not classify continuous source-core forms,
derive a Fourier-support law for them, or study their \(L^2\) radical and closability.

The two harmonic-analysis inputs are classical.  Bochner--Schwartz represents positive-definite
tempered distributions by positive temperate measures \cite{Schwartz1966}; Wiener's theorem
characterizes cyclic translation vectors by almost-everywhere nonvanishing of the Fourier
transform \cite{Wiener1932}.  Here the theta identity \(\widehat k=\Xi\) activates both theorems
at once.  Theta nullity forces the Bochner--Schwartz measure onto the real zeta divisor, while the
same divisor is Lebesgue-null and therefore makes \(k\) cyclic in ordinary \(L^2\).  These two
facts point in opposite topological directions: the admissible metric lives only on the divisor,
but the metric radical is dense in the ambient \(L^2\) representation.

General Hilbert--Pólya programs, de Branges spaces, trace-formula constructions, and
noncommutative-geometric approaches use different operators and domains
\cite{deBranges1968,Lagarias2006,Connes2026}.  The no-go theorem is not asserted for those
architectures.  It is source-specific to the Hedenmalm theta pair and exact for the regularity
class stated in Theorem \ref{thm:main}.

To our knowledge, the current literature surrounding the June 2026 theta proposal contains no earlier
divisor-support classification, dense-radical theorem, or \(L^2\)-nonclosability result for this
pair, nor the resulting positive-definite source-core no-go theorem.  The direct-parent comparison
above is the novelty boundary: classical representation and cyclicity theorems are used without
novelty claims.

\section{Conclusion}

The Hedenmalm theta pair admits a complete source-core metric classification.  Positivity and pair
symmetry force a unique positive temperate Fourier measure; theta nullity forces that measure to be
purely atomic on the real \(\Xi\)-divisor.  The same theta source is cyclic in ordinary \(L^2\), so
every nonzero classified metric has a dense proper nonclosed radical and a fixed-energy
\(L^2\)-null sequence.  Sequential \(L^2\)-closability occurs only for the zero form.

This closes the source-core Hilbert-metric problem for the new theta spectral program:
no positive-definite metric exists on \(L\cS(\R)\) under the natural continuity and pair-symmetry
hypotheses, and no nonzero semidefinite source-core metric can arise from a closed nonnegative
quadratic form on \(L^2(\R)\).  The exact replacement is not another local operator but a
degenerate divisor-supported Fourier rigging.  On theta eigenprofiles its conservativity criterion is
precisely RH plus simplicity and positive mass at every divisor point.  These statements delimit
the theorem sharply: they do not claim RH, and they do not extend the no-go beyond the
source-core category proved here.

\begin{thebibliography}{99}

\bibitem{Connes2026}
A.~Connes,
\emph{The Riemann Hypothesis: Past, Present and a Letter Through Time},
arXiv:2602.04022, 2026.

\bibitem{deBranges1968}
L.~de Branges,
\emph{Hilbert Spaces of Entire Functions},
Prentice--Hall, Englewood Cliffs, NJ, 1968.

\bibitem{Hedenmalm2026}
H.~Hedenmalm,
\emph{Spectral interpretation of Riemann zeta zeros},
arXiv:2606.17494, 2026.

\bibitem{Lagarias2006}
J.~C. Lagarias,
The Riemann hypothesis: arithmetic and geometry,
in \emph{Surveys in Noncommutative Geometry},
Clay Math. Proc. \textbf{6} (2006), 127--139.

\bibitem{Schwartz1966}
L.~Schwartz,
\emph{Th\'eorie des distributions},
new edition, Hermann, Paris, 1966.

\bibitem{Wiener1932}
N.~Wiener,
Tauberian theorems,
\emph{Ann. of Math.} (2) \textbf{33} (1932), 1--100.

\end{thebibliography}

\end{document}
